% `tex.web` §298 `print_cmd_chr`: a primitive is printed by *its own* name.
%
% This engine keeps the name a primitive was registered under in a side table,
% because deriving it from the Rust variant is wrong wherever the two differ —
% `\pdfshellescape` would print `\shellescape` and `\-` would print
% `\dischyphen`. Two entries were missing, and they are the two whose variant
% carries a suffix to dodge a name clash: `\meaning\string` said `\stringp` and
% `\meaning\meaning` said `\meaningp`. §367's level-2 command trace reads the
% same table, so both were wrong there too.
%
% `\meaning` through a `\let` alias is the sharper form of the same question:
% TeX prints the primitive it came from, never the name it was reached by.
\input plain
\let\zs=\string \let\zm=\meaning
\message{<\meaning\string><\meaning\meaning>}
\message{<\meaning\zs><\meaning\zm>}
\message{<\meaning\noexpand><\meaning\number><\meaning\romannumeral>}
\message{<\meaning\jobname><\meaning\fontname><\meaning\the>}
\message{<\meaning\csname><\meaning\endcsname><\meaning\expandafter>}
% ...and the one-character ones, which is why the table exists at all.
\message{<\meaning\-><\meaning\/><\meaning\ >}
\bye
